\chapter{Global Phase, Polarization Ellipses, and Unitary Evolution}

\section{Completing the Polarization Observable}

We begin with the unfinished exercise from the preceding lecture. The
horizontal and vertical polarization states are
\begin{equation}
  \ket{x}=
  \begin{pmatrix}1\\0\end{pmatrix},
  \qquad
  \ket{y}=
  \begin{pmatrix}0\\1\end{pmatrix},
\end{equation}
and the observable that records \(+1\) for horizontal transmission and
\(-1\) for vertical transmission is
\begin{equation}
  \hat P_0=
  \begin{pmatrix}
    1&0\\
    0&-1
  \end{pmatrix}.
  \label{eq:l7-P0}
\end{equation}
At \(45^\circ\), the two eigenstates
\begin{equation}
  \ket{d_\pm}
  =\frac{1}{\sqrt2}
    \begin{pmatrix}1\\ \pm1\end{pmatrix}
\end{equation}
belong to
\begin{equation}
  \hat P_{\pi/4}
  =
  \begin{pmatrix}
    0&1\\
    1&0
  \end{pmatrix}.
  \label{eq:l7-P45}
\end{equation}

The homework was to construct the observable for the circular basis
\begin{equation}
  \ket{c_+}
  =\frac{1}{\sqrt2}
    \begin{pmatrix}1\\i\end{pmatrix},
  \qquad
  \ket{c_-}
  =\frac{1}{\sqrt2}
    \begin{pmatrix}1\\-i\end{pmatrix}.
  \label{eq:l7-circular-basis}
\end{equation}
The spectral rule gives the answer immediately:
\begin{align}
  \hat P_c
  &=
  \ket{c_+}\bra{c_+}
  -\ket{c_-}\bra{c_-}\\
  &=
  \begin{pmatrix}
    0&-i\\
    i&0
  \end{pmatrix}.
  \label{eq:l7-Pc}
\end{align}
Indeed,
\begin{equation}
  \hat P_c\ket{c_\pm}=\pm\ket{c_\pm}.
\end{equation}
The only delicate point in the check is \(i^2=-1\), together with complex
conjugation when a ket becomes a bra.

\begin{figure}[tbp]
  \centering
  \includegraphics[width=0.70\textwidth]{lecture_07_figure_02.png}
  \caption{The three polarization bases and the circular observable at the
  end of the homework review.}
  \lectureframe{00:05:20}
\end{figure}

Orthogonality has an operational meaning throughout this discussion. Two
orthogonal polarization states can be distinguished with certainty in one
ideal analyzer: an apparatus that transmits one rejects the other. This is
why the inner product is not decorative notation. It records a measurable
relation between preparations.

\section{What States and Observables Are For}

It is useful to pause and compare our present position with the opening
classical discussion. For a coin or die, the state is one member of a finite
set. A classical observable is simply a numerical function on that set. We
might assign \(+1\) to heads and \(-1\) to tails, while the die arrives with
the labels \(1,\ldots,6\). Relabelling the outcomes changes the numerical
convention, not the underlying states.

The same idea survives in continuous classical mechanics. A state is a point
\((q,p)\) in phase space, and an observable is a function
\begin{equation}
  F=F(q,p).
\end{equation}
Position \(q\), momentum \(p\), and \(p+q\) are elementary examples. As the
phase-space point moves, the numerical value of each function changes.

Quantum mechanics replaces the set of states by a complex vector space and
the numerical functions by Hermitian operators. Before discussing motion in
that space, however, we should isolate the two principal questions that a
state vector already answers.

First, for a normalized state \(\ket{\psi}\) and observable \(\hat K\), the
ensemble average is
\begin{equation}
  \expect{\hat K}_\psi
  =\bra{\psi}\hat K\ket{\psi}.
  \label{eq:l7-expectation}
\end{equation}
Prepare the same state repeatedly, measure \(K\) on each member of the
ensemble, and the empirical mean approaches this number.

Second, suppose
\begin{equation}
  \hat K\ket{a}=\alpha\ket{a},
  \qquad
  \hat L\ket{b}=\beta\ket{b}.
\end{equation}
If the preparation makes \(K=\alpha\) definite and the subsequent
measurement asks whether \(L=\beta\), then
\begin{equation}
  \operatorname{Prob}(b|a)
  =|\braket{b}{a}|^2
  =\braket{a}{b}\braket{b}{a}.
  \label{eq:l7-transition-probability}
\end{equation}
Reversing the preparation and test gives the same numerical probability,
because \(|\braket{a}{b}|^2=|\braket{b}{a}|^2\).

\begin{figure}[tbp]
  \centering
  \includegraphics[width=0.74\textwidth]{lecture_07_figure_03.png}
  \caption{Expectation value, eigenstate preparations, and the squared
  overlap as the two physical uses of a state vector.}
  \lectureframe{00:14:20}
\end{figure}

\begin{classroomqa}
\audiencequestion{Does the order of the two overlaps in the probability
matter?\lecturetimestamp{00:14:54--00:15:18}}
\lecturerresponse{The two amplitudes are complex conjugates:
\(\braket{a}{b}=\braket{b}{a}^{*}\). Their product is therefore the same
nonnegative real number in either order. The correction on the board changes
the amplitude, but not the transition probability.}
\end{classroomqa}

Average values and transition probabilities both refer to repeated
experiments. A single trial returns one allowed outcome; the probability and
the average emerge from an ensemble. This common structure leads directly to
the question of phase.

\section{A Physical State Is a Ray}

\subsection{Pure Phases}

A physical state vector is normalized:
\begin{equation}
  \braket{\psi}{\psi}=1.
  \label{eq:l7-normalization}
\end{equation}
If \(\ket{\psi}=\sum_n c_n\ket{n}\) in an orthonormal basis, this is
\(\sum_n|c_n|^2=1\), the statement that the total probability is one.

Normalization does not remove every redundancy. To see what remains, write a
complex number in polar form,
\begin{equation}
  z=x+iy=re^{i\chi}
  =r(\cos\chi+i\sin\chi).
\end{equation}
Multiplication multiplies moduli and adds angles:
\begin{equation}
  r_1e^{i\chi_1}r_2e^{i\chi_2}
  =r_1r_2e^{i(\chi_1+\chi_2)}.
\end{equation}
Complex conjugation reflects the number across the real axis,
\begin{equation}
  z^*=re^{-i\chi}.
\end{equation}
A complex number of unit modulus is a \emph{pure phase}:
\begin{equation}
  e^{i\chi},
  \qquad
  e^{i\chi}e^{-i\chi}=1.
  \label{eq:l7-pure-phase}
\end{equation}
The word phase may refer either to the angle \(\chi\) or to the unit complex
number \(e^{i\chi}\); the context tells us which is intended.

\begin{figure}[tbp]
  \centering
  \includegraphics[width=0.68\textwidth]{lecture_07_figure_04.png}
  \caption{The unit circle beside the comparison of
  \(\ket{A}\) and \(e^{i\chi}\ket{A}\).}
  \lectureframe{00:21:40}
\end{figure}

\subsection{Why the Overall Phase Cancels}

Mathematically, \(\ket{\psi}\) and \(e^{i\chi}\ket{\psi}\) are different
vectors. Physically, they represent the same pure state. Let
\begin{equation}
  \ket{\psi'}=e^{i\chi}\ket{\psi}.
\end{equation}
The corresponding bra is
\begin{equation}
  \bra{\psi'}=e^{-i\chi}\bra{\psi}.
\end{equation}
Every expectation value is unchanged:
\begin{align}
  \bra{\psi'}\hat K\ket{\psi'}
  &=
  e^{-i\chi}\bra{\psi}\hat K
  e^{i\chi}\ket{\psi}\\
  &=
  \bra{\psi}\hat K\ket{\psi}.
  \label{eq:l7-phase-expectation}
\end{align}
The phase is an ordinary scalar and commutes through \(\hat K\).

The transition probability is equally insensitive. If
\(\ket{b'}=e^{i\eta}\ket{b}\), then
\begin{align}
  |\braket{b'}{a}|^2
  &=
  e^{-i\eta}\braket{b}{a}
  e^{i\eta}\braket{a}{b}\\
  &=
  |\braket{b}{a}|^2.
  \label{eq:l7-phase-probability}
\end{align}
Independent phase choices for either state disappear from the absolute
square.

\begin{figure}[tbp]
  \centering
  \includegraphics[width=0.72\textwidth]{lecture_07_figure_05.png}
  \caption{The two phase factors cancel in both the operator sandwich and
  the squared overlap.}
  \lectureframe{00:23:35}
\end{figure}

The physical object is therefore not one normalized vector but the
equivalence class
\begin{equation}
  \ket{\psi}\sim e^{i\chi}\ket{\psi}.
  \label{eq:l7-ray-equivalence}
\end{equation}
Such an equivalence class is called a \emph{ray} in Hilbert space. In
particular,
\begin{equation}
  \begin{pmatrix}0\\1\end{pmatrix}
  \sim
  \begin{pmatrix}0\\-1\end{pmatrix},
\end{equation}
because \(-1=e^{i\pi}\). Likewise,
\begin{equation}
  \frac{1}{\sqrt2}
  \begin{pmatrix}1\\i\end{pmatrix}
  \sim
  \frac{1}{\sqrt2}
  \begin{pmatrix}i\\-1\end{pmatrix},
  \label{eq:l7-circular-phase-example}
\end{equation}
because the second vector is \(i\) times the first.

\begin{figure}[tbp]
  \centering
  \includegraphics[width=0.72\textwidth]{lecture_07_figure_06.png}
  \caption{Multiplying both circular-polarization components by \(i\)
  changes the vector but not the physical ray.}
  \lectureframe{00:28:40}
\end{figure}

\begin{classroomqa}
\audiencequestion{What notation should express that two vectors have the
same physical significance without saying that they are equal?
\lecturetimestamp{00:25:49--00:27:50}}
\lecturerresponse{Ordinary equality is too strong, because the vectors
really differ. The standard statement is that they represent the same ray,
written here as \(\ket{\psi}\sim e^{i\chi}\ket{\psi}\). What is preserved is
every expectation value and transition probability.}
\end{classroomqa}

\section{Two Parameters Describe Polarization}

Begin with the general two-component vector
\begin{equation}
  \ket{\psi}
  =
  \begin{pmatrix}\alpha\\\beta\end{pmatrix},
  \qquad
  \alpha,\beta\in\mathbb C.
  \label{eq:l7-general-spinor}
\end{equation}
The two complex entries contain four real parameters. Normalization removes
one:
\begin{equation}
  |\alpha|^2+|\beta|^2=1.
\end{equation}
The irrelevant common phase removes another. Except at the coordinate pole
\(\alpha=0\), choose that phase so that \(\alpha=a\) is real and
nonnegative. Then
\begin{equation}
  \beta=\sqrt{1-a^2}\,e^{i\phi},
  \qquad
  0\leq a\leq1,
\end{equation}
and
\begin{equation}
  \ket{\psi}
  =
  \begin{pmatrix}
    a\\
    \sqrt{1-a^2}\,e^{i\phi}
  \end{pmatrix}.
  \label{eq:l7-gauge-fixed-state}
\end{equation}
Exactly two real parameters remain: \(a\) and the relative phase \(\phi\).

\begin{figure}[tbp]
  \centering
  \includegraphics[width=0.68\textwidth]{lecture_07_figure_07.png}
  \caption{Normalization and phase fixing reduce two complex amplitudes to
  two physical real parameters.}
  \lectureframe{00:32:45}
\end{figure}

An equivalent parameterization makes the geometry especially transparent:
\begin{equation}
  \ket{\psi}
  =
  \begin{pmatrix}
    \cos(\vartheta/2)\\
    e^{i\phi}\sin(\vartheta/2)
  \end{pmatrix},
  \qquad
  0\leq\vartheta\leq\pi.
  \label{eq:l7-bloch-parameterization}
\end{equation}
The rays of a two-state system form a sphere, often called the Bloch sphere
or, for polarization, the Poincare sphere. The lecture does not require that
name, but its two-parameter count is precisely this geometry.

\section{From the State Vector to a Polarization Ellipse}

For a monochromatic wave moving along \(z\), write the transverse electric
field as
\begin{equation}
  \bm E(z,t)
  =
  \operatorname{Re}
  \left[
    E_0
    \left(
      \alpha\,\hat{\bm x}
      +\beta\,\hat{\bm y}
    \right)
    e^{i(kz-\omega t)}
  \right].
  \label{eq:l7-Jones-field}
\end{equation}
The normalized pair \((\alpha,\beta)^T\) is the Jones vector. If the two
components have the same phase, the tip of \(\bm E\) moves back and forth
along a line. A quarter-cycle relative phase with equal magnitudes produces
a circle. General amplitudes and relative phases produce an ellipse.

The examples on the board may be written dimensionally as
\begin{align}
  E_x&=E_0\cos(kz-\omega t),\\
  E_y&=E_0\sin(kz-\omega t),
  \label{eq:l7-circle-field}
\end{align}
which traces a circle at fixed \(z\), and
\begin{align}
  E_x&=E_0\cos(kz-\omega t),\\
  E_y&=\frac{E_0}{2}\sin(kz-\omega t),
  \label{eq:l7-ellipse-field}
\end{align}
which traces an ellipse with its major axis along \(x\).

\begin{figure}[tbp]
  \centering
  \includegraphics[width=0.70\textwidth]{lecture_07_figure_08.png}
  \caption{Equal quadrature amplitudes give a circle; reducing one amplitude
  gives an ellipse.}
  \lectureframe{00:35:25}
\end{figure}

The ellipse needs two independent real parameters. One gives the orientation
of its major axis. The other gives its shape together with the sense of
rotation. The lecture informally calls the latter a signed eccentricity. In
standard optics it is cleaner to use the \emph{signed ellipticity angle}
\(\epsilon\):
\begin{equation}
  -\frac{\pi}{4}\leq\epsilon\leq\frac{\pi}{4},
  \qquad
  \left|\tan\epsilon\right|
  =\frac{\text{minor semiaxis}}{\text{major semiaxis}}.
  \label{eq:l7-ellipticity-angle}
\end{equation}
The sign records handedness. Linear polarization has \(\epsilon=0\);
circular polarization has \(|\epsilon|=\pi/4\), where the major-axis
direction becomes meaningless.

\begin{classroomqa}
\audiencequestion{Do clockwise and counterclockwise polarization require a
third parameter?\lecturetimestamp{00:37:52--00:39:00}}
\lecturerresponse{No. The shape parameter may carry a sign. As an ellipse is
squeezed through the linear-polarization limit, changing the sign reverses
the sense in which the field traverses the curve. Standard optics expresses
this with signed ellipticity rather than literal signed geometric
eccentricity.}
\end{classroomqa}

\subsection{Reading the Ellipse from Analyzer Probabilities}

Let
\begin{equation}
  b=\sqrt{1-a^2},
  \qquad
  \ket{\psi}
  =
  \begin{pmatrix}a\\be^{i\phi}\end{pmatrix},
\end{equation}
and compare it with a linearly polarized analyzer state
\begin{equation}
  \ket{\theta}
  =
  \begin{pmatrix}
    \cos\theta\\
    \sin\theta
  \end{pmatrix}.
\end{equation}
The two conjugate amplitudes are
\begin{align}
  \braket{\theta}{\psi}
  &=a\cos\theta+be^{i\phi}\sin\theta,\\
  \braket{\psi}{\theta}
  &=a\cos\theta+be^{-i\phi}\sin\theta.
  \label{eq:l7-ellipse-overlaps}
\end{align}

\begin{classroomqa}
\audiencequestion{Should one of the phase factors in the overlap be complex
conjugated?\lecturetimestamp{00:41:48--00:42:02}}
\lecturerresponse{Yes. The sign depends on the bra-ket order. In
\(\braket{\psi}{\theta}\), the coefficient from \(\ket{\psi}\) is
conjugated, giving \(e^{-i\phi}\). The opposite order gives \(e^{i\phi}\).
Their absolute squares agree.}
\end{classroomqa}

The transmission probability is
\begin{align}
  P_\theta
  &=
  |\braket{\theta}{\psi}|^2\\
  &=
  a^2\cos^2\theta
  +b^2\sin^2\theta
  +2ab\cos\phi\,\sin\theta\cos\theta\\
  &=
  \frac12
  \left[
    1
    +(a^2-b^2)\cos2\theta
    +2ab\cos\phi\,\sin2\theta
  \right].
  \label{eq:l7-linear-analyzer-probability}
\end{align}
The lecture leaves this multiplication and maximization as homework. Finishing
that step makes its geometric prescription exact. The major-axis angle
\(\theta_{\mathrm{maj}}\) obeys
\begin{equation}
  2\theta_{\mathrm{maj}}
  =
  \operatorname{atan2}
  \left(2ab\cos\phi,\ a^2-b^2\right)
  \pmod{2\pi}.
  \label{eq:l7-major-axis}
\end{equation}
Thus the analyzer most likely to transmit the photon points along the major
axis of the polarization ellipse.

\begin{figure}[tbp]
  \centering
  \includegraphics[width=0.78\textwidth]{lecture_07_figure_09.png}
  \caption{The overlap amplitude and the geometric search for the
  linear-analyzer direction of maximum transmission.}
  \lectureframe{00:45:20}
\end{figure}

It is useful to collect the same information in the normalized Stokes
parameters
\begin{equation}
  S_1=a^2-b^2,
  \qquad
  S_2=2ab\cos\phi,
  \qquad
  S_3=2ab\sin\phi,
  \label{eq:l7-stokes}
\end{equation}
for which \(S_1^2+S_2^2+S_3^2=1\). With a fixed propagation and time
convention,
\begin{equation}
  2\theta_{\mathrm{maj}}=\operatorname{atan2}(S_2,S_1),
  \qquad
  \sin2\epsilon=S_3.
  \label{eq:l7-orientation-ellipticity}
\end{equation}
Changing the convention for viewing the wave reverses the sign assigned to
handedness, but not the two-parameter count.

\begin{figure}[tbp]
  \centering
  \includegraphics[width=0.68\textwidth]{lecture_07_figure_10.png}
  \caption{Squeezing through the linear limit reverses the traversal sense,
  which is encoded by the sign of ellipticity.}
  \lectureframe{00:46:40}
\end{figure}

\begin{classroomqa}
\audiencequestion{Orientation and ordinary eccentricity still seem to omit
clockwise versus counterclockwise motion. Where is it stored?
\lecturetimestamp{00:45:31--00:47:14}}
\lecturerresponse{It is stored in the sign of the second parameter. The
major-axis angle gives the orientation; signed ellipticity gives both the
axis ratio and the traversal sense. At the linear limit the sign can change,
and at the circular limit the axis angle is irrelevant.}
\end{classroomqa}

\section{Time Evolution Must Preserve Distinctions}

\subsection{The Classical Analogy}

We now turn from the description of states to their motion. A finite
classical state set admits only jumps if time is viewed microscopically, but
the quantum vector space contains continuous paths even when a measurement
has only two possible outcomes. A linear polarization axis can rotate
continuously with time although every fixed analyzer still records only
\(+1\) or \(-1\).

For a closed reversible classical system, evolution is one-to-one. It may
permute the states, but it cannot send two distinct initial states to the
same final state. If both \(5\) and \(2\) on a die evolved to \(3\), the final
state would not reveal which initial state had been prepared. The distinction
would have been erased.

\begin{figure}[tbp]
  \centering
  \includegraphics[width=0.66\textwidth]{lecture_07_figure_11.png}
  \caption{A forbidden many-to-one rule would send both \(5\) and \(2\) to
  \(3\), erasing the initial distinction.}
  \lectureframe{00:51:15}
\end{figure}

Quantum states have more structure than simply same or different. For
normalized vectors, the invariant transition probability
\begin{equation}
  |\braket{B}{A}|^2
  \label{eq:l7-ray-overlap}
\end{equation}
measures their operational overlap. Orthogonal states have zero overlap and
can be distinguished with certainty. States on the same ray have
\(|\braket{B}{A}|=1\). Intermediate values mean that no single ideal test
distinguishes the preparations with certainty.

\begin{figure}[tbp]
  \centering
  \includegraphics[width=0.62\textwidth]{lecture_07_figure_12.png}
  \caption{Zero overlap represents orthogonal, perfectly distinguishable
  polarization states.}
  \lectureframe{00:54:40}
\end{figure}

\subsection{Preserving the Inner Product}

Prepare two identical systems in \(\ket{A}\) and \(\ket{B}\), and let both
evolve for the same time:
\begin{equation}
  \ket{A}\longrightarrow\ket{A'},
  \qquad
  \ket{B}\longrightarrow\ket{B'}.
\end{equation}
Closed-system quantum evolution preserves their inner product:
\begin{equation}
  \boxed{\braket{B'}{A'}=\braket{B}{A}.}
  \label{eq:l7-inner-product-preservation}
\end{equation}
Consequently equal rays remain equal rays, orthogonal states remain
orthogonal, and every transition probability remains unchanged. This is the
quantum form of preserving distinctions.

\begin{figure}[tbp]
  \centering
  \includegraphics[width=0.72\textwidth]{lecture_07_figure_13.png}
  \caption{The overlap of two states is unchanged when both undergo the same
  closed-system evolution.}
  \lectureframe{00:57:40}
\end{figure}

\begin{classroomqa}
\audiencequestion{If one photon fails a horizontal analyzer, how do we know
it was vertical rather than a \(45^\circ\) state that happened to fail?
\lecturetimestamp{00:58:17--00:59:00}}
\lecturerresponse{From one event, we do not. Orthogonality is tested by
certainty across identically prepared trials. Every genuinely vertical
photon is rejected by an ideal horizontal analyzer; a \(45^\circ\) state is
rejected only with probability one half. The ensemble distinguishes the
preparations.}
\end{classroomqa}

\begin{classroomqa}
\audiencequestion{Does the uncertainty principle force a small chance that a
perfectly horizontal photon passes a vertical polarizer?
\lecturetimestamp{00:59:01--01:00:31}}
\lecturerresponse{No. One polarization observable can be prepared with
arbitrary precision in the ideal theory. The uncertainty concerns
incompatible observables: no state has definite outcomes for two genuinely
different analyzer axes at once. It does not blur an eigenstate of one chosen
axis.}
\end{classroomqa}

In an ordinary real vector space, a transformation preserving all inner
products looks like a rigid rotation or reflection: vectors move, but their
lengths and mutual angles do not. The complex quantum version is a unitary
transformation.

\begin{figure}[tbp]
  \centering
  \includegraphics[width=0.60\textwidth]{lecture_07_figure_14.png}
  \caption{The geometric analogy: vectors move while their lengths and
  mutual angles remain fixed.}
  \lectureframe{01:01:25}
\end{figure}

\section{Operators on Bras and the Adjoint}

Let \(\hat L\) be a linear operator on kets:
\begin{equation}
  \hat L\ket{A}=\ket{C}.
  \label{eq:l7-LA-C}
\end{equation}
Linearity means
\begin{equation}
  \hat L(\alpha\ket{A}+\beta\ket{B})
  =
  \alpha\hat L\ket{A}
  +\beta\hat L\ket{B}.
\end{equation}
For any bra \(\bra{B}\), the matrix element
\begin{equation}
  \bra{B}\hat L\ket{A}
\end{equation}
may be read associatively as the inner product of \(\bra{B}\) with
\(\hat L\ket{A}\), or as the bra \(\bra{B}\hat L\) applied to
\(\ket{A}\).

Inserting a complete orthonormal basis makes the latter bra explicit:
\begin{equation}
  \bra{B}\hat L
  =
  \sum_n
  \bra{B}\hat L\ket{n}\bra{n}.
  \label{eq:l7-operator-on-bra}
\end{equation}
The coefficients are numbers, so the right side is a linear combination of
bras.

\begin{figure}[tbp]
  \centering
  \includegraphics[width=0.72\textwidth]{lecture_07_figure_15.png}
  \caption{A completeness insertion defines the bra
  \(\bra{B}\hat L\) from the known action of \(\hat L\) on kets.}
  \lectureframe{01:06:50}
\end{figure}

\begin{classroomqa}
\audiencequestion{Does this construction assume that \(\hat L\) is
Hermitian?\lecturetimestamp{01:05:51--01:06:00}}
\lecturerresponse{No. The matrix element and the bra
\(\bra{B}\hat L\) are defined for a general linear operator. Hermiticity is
the additional condition \(\hat L^\dagger=\hat L\), introduced later.}
\end{classroomqa}

\begin{classroomqa}
\audiencequestion{Are we assuming associativity when the brackets are
regrouped?\lecturetimestamp{01:06:01--01:06:18}}
\lecturerresponse{Yes. Operator composition and its action on vectors are
associative wherever the products are defined:
\(\bra{B}(\hat L\ket{A})=(\bra{B}\hat L)\ket{A}\). This is what lets the
notation be manipulated consistently.}
\end{classroomqa}

\subsection{Hermitian Conjugation}

Taking the adjoint of Equation~\eqref{eq:l7-LA-C} gives
\begin{equation}
  \bra{A}\hat L^\dagger=\bra{C}.
  \label{eq:l7-adjoint-map}
\end{equation}
The defining matrix-element identity is
\begin{equation}
  \boxed{
  \bra{B}\hat L\ket{A}
  =
  \left(
    \bra{A}\hat L^\dagger\ket{B}
  \right)^*.}
  \label{eq:l7-adjoint-definition}
\end{equation}
It combines two operations already familiar from inner products:
interchange bra and ket, then complex conjugate.

\begin{figure}[tbp]
  \centering
  \includegraphics[width=0.74\textwidth]{lecture_07_figure_16.png}
  \caption{The adjoint relation obtained by reversing the bra-ket order and
  complex conjugating the matrix element.}
  \lectureframe{01:12:50}
\end{figure}

In an orthonormal basis, define
\begin{equation}
  L_{mn}=\bra{m}\hat L\ket{n}.
\end{equation}
Then
\begin{equation}
  (\hat L^\dagger)_{mn}=L_{nm}^*.
  \label{eq:l7-dagger-components}
\end{equation}
Thus the matrix of \(\hat L^\dagger\) is the complex-conjugate transpose:
\begin{equation}
  L^\dagger=(L^*)^T.
\end{equation}
Rows and columns are interchanged, and every entry is conjugated.

\begin{figure}[tbp]
  \centering
  \includegraphics[width=0.72\textwidth]{lecture_07_figure_17.png}
  \caption{The matrix rule for the adjoint: transpose and complex
  conjugate.}
  \lectureframe{01:17:25}
\end{figure}

\begin{classroomqa}
\audiencequestion{Are we assuming a finite-dimensional state space here?
\lecturetimestamp{01:13:44--01:14:05}}
\lecturerresponse{For the displayed matrices, yes. The same adjoint identity
extends to Hilbert spaces, but infinite-dimensional operators require care
about domains and completeness. The finite-dimensional case contains the
algebra needed for the present argument without those analytic
qualifications.}
\end{classroomqa}

\section{Hermitian Operators Represent Observables}

An operator is Hermitian when
\begin{equation}
  \hat L=\hat L^\dagger.
  \label{eq:l7-Hermitian}
\end{equation}
In components,
\begin{equation}
  L_{mn}=L_{nm}^*.
\end{equation}
The diagonal entries are therefore real,
\begin{equation}
  L_{nn}=L_{nn}^*\in\mathbb R,
\end{equation}
while entries reflected across the diagonal are complex conjugates.

\begin{figure}[tbp]
  \centering
  \includegraphics[width=0.68\textwidth]{lecture_07_figure_18.png}
  \caption{The conjugate-symmetry condition
  \(L_{mn}=L_{nm}^{*}\) for a Hermitian matrix.}
  \lectureframe{01:19:30}
\end{figure}

In finite dimensions, the spectral theorem gives the properties needed of
an observable:
\begin{enumerate}
  \item every eigenvalue is real;
  \item eigenvectors belonging to distinct eigenvalues are orthogonal;
  \item an orthonormal eigenbasis exists.
\end{enumerate}
The eigenvectors need not have real components. Circular polarization is the
immediate counterexample: the Hermitian matrix \(\hat P_c\) has eigenvectors
containing \(\pm i\).

The expectation value of a Hermitian operator is real:
\begin{align}
  \left(
    \bra{\psi}\hat L\ket{\psi}
  \right)^*
  &=
  \bra{\psi}\hat L^\dagger\ket{\psi}\\
  &=
  \bra{\psi}\hat L\ket{\psi}.
\end{align}
Conversely, in finite dimensions, if
\(\bra{\psi}\hat L\ket{\psi}\) is real for every \(\ket{\psi}\), then
\(\hat L\) is Hermitian.

\begin{classroomqa}
\audiencequestion{How large is the matrix representing an observable?
\lecturetimestamp{01:20:28--01:22:28}}
\lecturerresponse{It matches the dimension of the state space under study.
Polarization alone is two-dimensional. A full photon state also carries
continuous momentum or position labels and is generally
infinite-dimensional. Finite truncations are often used for calculation,
but finite systems such as an isolated spin degree of freedom are not merely
approximations.}
\end{classroomqa}

The three polarization observables make a compact check:
\begin{equation}
  \hat P_0=
  \begin{pmatrix}1&0\\0&-1\end{pmatrix},
  \quad
  \hat P_{\pi/4}=
  \begin{pmatrix}0&1\\1&0\end{pmatrix},
  \quad
  \hat P_c=
  \begin{pmatrix}0&-i\\i&0\end{pmatrix}.
  \label{eq:l7-three-Pauli}
\end{equation}
Each equals its adjoint. They are, respectively, the Pauli matrices
\(\sigma_z,\sigma_x,\sigma_y\).

\begin{figure}[tbp]
  \centering
  \includegraphics[width=0.68\textwidth]{lecture_07_figure_19.png}
  \caption{The axial, diagonal, and circular polarization observables are
  all Hermitian.}
  \lectureframe{01:24:45}
\end{figure}

\begin{classroomqa}
\audiencequestion{Why is Hermiticity the right condition for an observable?
\lecturetimestamp{01:25:08--01:26:20}}
\lecturerresponse{Laboratory outcomes and their averages are real numbers.
Hermiticity guarantees real eigenvalues and real expectation values. In
finite dimensions the converse condition on all quadratic forms also
recovers Hermiticity, so the mathematical requirement is tightly matched to
the operational one.}
\end{classroomqa}

\begin{classroomqa}
\audiencequestion{Must the eigenvectors of an observable span the whole
state space?\lecturetimestamp{01:26:22--01:27:43}}
\lecturerresponse{Yes in the finite-dimensional Hermitian case: the spectral
theorem supplies a complete orthonormal eigenbasis, including an orthonormal
choice inside degenerate eigenspaces. A general non-Hermitian operator may
have nonorthogonal eigenvectors or may fail to have a complete eigenbasis.}
\end{classroomqa}

\section{Unitary Operators Preserve Quantum Relations}

Hermitian operators represent quantities that can be measured. They are not,
in general, the transformations that move states while preserving their
relations. Let one linear operator \(\hat U\) act on every state:
\begin{equation}
  \hat U\ket{A}=\ket{A'},
  \qquad
  \hat U\ket{B}=\ket{B'}.
  \label{eq:l7-U-map}
\end{equation}
The transformed bras are
\begin{equation}
  \bra{A'}=\bra{A}\hat U^\dagger,
  \qquad
  \bra{B'}=\bra{B}\hat U^\dagger.
\end{equation}
Therefore
\begin{equation}
  \braket{B'}{A'}
  =
  \bra{B}\hat U^\dagger\hat U\ket{A}.
\end{equation}
Demanding equality with \(\braket{B}{A}\) for every pair
\(\ket{A},\ket{B}\) forces
\begin{equation}
  \boxed{\hat U^\dagger\hat U=\mathbf 1.}
  \label{eq:l7-isometry}
\end{equation}

In finite dimensions this condition makes \(\hat U\) invertible and implies
\begin{equation}
  \hat U^\dagger=\hat U^{-1},
  \qquad
  \hat U\hat U^\dagger
  =\hat U^\dagger\hat U
  =\mathbf 1.
  \label{eq:l7-unitary}
\end{equation}
In an infinite-dimensional Hilbert space, the definition of a unitary
operator includes both equalities; \(\hat U^\dagger\hat U=\mathbf1\) alone
defines an isometry that need not be onto.

\begin{figure}[tbp]
  \centering
  \includegraphics[width=0.74\textwidth]{lecture_07_figure_20.png}
  \caption{Preservation of every inner product leads to the unitary
  condition \(\hat U^\dagger\hat U=\mathbf1\).}
  \lectureframe{01:34:50}
\end{figure}

A unitary transformation is not the identity. It may rotate polarization,
implement a symmetry, or carry a state from one time to another. Its adjoint
undoes it, and throughout the motion
\begin{equation}
  \braket{UB}{UA}
  =
  \bra{B}\hat U^\dagger\hat U\ket{A}
  =
  \braket{B}{A}.
\end{equation}
Thus every norm, orthogonality relation, and transition probability is
preserved.

\begin{classroomqa}
\audiencequestion{Does preserving relationships mean anything broader than
preserving inner products?\lecturetimestamp{01:35:13--01:35:53}}
\lecturerresponse{Mathematically, inner-product preservation is the precise
condition used here. Physically it preserves the absolute squares of those
inner products and therefore every transition probability between
simultaneously evolved preparations.}
\end{classroomqa}

The lecture closes just as the eigenvalue structure of a unitary operator is
introduced:
\begin{equation}
  \hat U\ket{n}=\lambda_n\ket{n}.
\end{equation}
The unanswered question has a short consequence of unitarity:
\begin{align}
  1
  &=
  \braket{n}{n}
  =
  \bra{n}\hat U^\dagger\hat U\ket{n}\\
  &=
  |\lambda_n|^2\braket{n}{n},
\end{align}
so
\begin{equation}
  |\lambda_n|=1,
  \qquad
  \lambda_n=e^{i\omega_n}.
  \label{eq:l7-unitary-eigenvalue}
\end{equation}
Finite-dimensional unitary matrices are normal and therefore possess a
complete orthonormal eigenbasis, but their eigenvalues are phases rather than
necessarily real numbers.

\begin{figure}[H]
  \centering
  \includegraphics[width=0.62\textwidth]{lecture_07_figure_21.png}
  \caption{The final board question: the eigenvalues of a unitary operator
  lie on the unit circle.}
  \lectureframe{01:37:20}
\end{figure}

We have reached the threshold of dynamics. Closed systems evolve by unitary
transformations. The next step is to discover how continuous unitary
evolution is generated by a Hermitian Hamiltonian and how that statement
becomes the Schrodinger equation.
